\documentclass[	pdftex, 
								a4paper,
								11pt, DIV11, BCOR5mm,
								parskip,
								%openany,
								]{scrreprt}
\usepackage[utf8]{inputenc}
\usepackage[T1]{fontenc}

\usepackage[colorlinks=true,linkcolor=blue,citecolor=blue]{hyperref}
\usepackage{comment}
\usepackage[ngerman]{babel}
\usepackage[nonumberlist,toc,acronym]{glossaries}

\usepackage{fancyhdr}
\usepackage[section]{placeins}
\usepackage{graphicx}
\usepackage{listings}
\usepackage[german]{babelbib}

\usepackage{varwidth}
\usepackage{float}

\usepackage{epstopdf}
\usepackage{textpos}
\usepackage{changepage}
\usepackage{theorem}
\usepackage{caption}
\usepackage{framed}
\usepackage{booktabs}
\usepackage{tikz}
\usetikzlibrary{decorations.markings,arrows,positioning,trees,calc,fit,shapes}
\usepackage[noend]{algorithmic}
\usepackage{algorithm}
\usepackage{environ}
\usepackage{ulem}

%Kopfzeile hinzufuegen
\renewcommand{\sectionmark}[1]{\markboth{#1}{}} % set the \leftmark

\fancypagestyle{plain}{
    
    \fancyhead[R]{\leftmark}
    \fancyhead[L]{\myauthor}
    \renewcommand{\headrulewidth}{1pt}
    %\fancyfoot[R]{\thepage}
}

\pagestyle{fancy}

%Literaturverzeichnis hinzufuegen
\bibliographystyle{babplain}

% Glossar hinzufuegen
\makeglossaries
\newacronym{usda} {USDA}{US Department of Agriculture}
\newacronym{fdc}{FDC}{FoodData Central}
\newacronym{rest}{REST}{Representational state transfer}
\newacronym{dge}{DGE}{Deutsche Gesellschaft für Ernährung}
\newacronym{pal}{PAL-Value}{physical activity level-Value}
\newacronym{json}{JSON}{JavaScript Object Notation}
\newacronym{iri}{IRI}{International Resource Identifier}
\newacronym{http}{HTTP}{HyperText Transfer Protocol}
\newacronym{xml}{XML}{Extensible Markup Language}
\newacronym{cdno}{CDNO}{Compositional Dietary Nutrition Ontology}
\newacronym{obo}{OBO}{Open Biological and Biomedical Ontology}
\newacronym{sparql}{SPARQL}{SPARQL Protocol And RDF Query Language}
\newacronym{sql}{SQL}{Structured Query Language}
\newacronym{uri}{URI}{Uniform Resource Identifier}
\newacronym{gi}{GI}{Glykämischer Index}
\newacronym{bmi}{BMI}{Body-Mass-Index}



\begin{document}
\setcounter{page}{-2}
\pagestyle{empty}
\pagenumbering{none}
\newcommand{\firstreviewer}{Prof. Michael Beetz PhD}
\newcommand{\secondreviewer}{}
\newcommand{\supervisor}{Michaela Kümpel}
\newcommand{\thesistype}{Master Thesis}
\newcommand{\myauthor}{Naser Azizi, Sorin Arion}
\newcommand{\mymaintitle}{A Knowledgebase with which you can generate robot plan for multiple mixing actions}
\newcommand{\mysubtitle}{A nutrition ontology with web interface for customer specific dietary \newline recommendations} 
\newcommand{\mytitle}{\centering {\Huge \mymaintitle}\\[.3in] 
    {\Large \mysubtitle}}
\newcommand{\pdftitle}{\mymaintitle - \mysubtitle}
\newcommand{\formattedfronttitle}{{\mytitle}}
\newcommand{\formattedinnertitle}{{\mytitle}}

\begin{titlepage}
	\vspace*{-2.2cm}
	\begin{adjustwidth}{-0cm}{-2.3cm}
	\thispagestyle{empty}
        \begin{figure}
            \begin{minipage}{.4\linewidth}
	\begin{flushleft}
		\includegraphics[height=1.5cm]{Graphics/unilogo-transp.pdf}
	\end{flushleft}
    \end{minipage}
    \hspace{.2\linewidth}
            \begin{minipage}{.4\linewidth}
	\begin{flushright}
		\includegraphics[height=3.0cm]{Graphics/logo-ai-small.pdf}
	\end{flushright}
    \end{minipage}
\end{figure}
	  \vfill
	  %\scalebox{0.95}{
    	  %\begin{minipage}{1.2\textwidth}
    	  %{
    		%\formattedfronttitle 
    	  %}
    	  %\end{minipage}
	  %}
	\begin{center}
	  {\huge \mymaintitle}\\
	  
	  \vfill
	  {\Large \thesistype}\\[2.5ex]
	  {\Large\em \myauthor}
	  \vfill
	{  
      \renewcommand\arraystretch{1.5}
      \begin{tabular}{l@{\hspace{2em}}r@{\hspace{1ex}}p{7cm}}
     Pr\"ufer der \thesistype: & 1. & \firstreviewer\\
                                 & 2. & \secondreviewer\\
	Supervisor		    &   & \supervisor\\
   \end{tabular}
  }
	\end{center}
	\end{adjustwidth}
\end{titlepage}
	\pagenumbering{arabic}
	\setcounter{page}{1}
	\pagestyle{plain}
	%\addcontentsline{toc}{chapter}{Eidesstaatliche Erkl\"arung}
	\chapter*{Eidesstattliche Erkl\"arung}
	
	Hiermit erkl\"aren wir, dass die vorliegende Arbeit selbstst\"andig angefertigt,
	nicht anderweitig zu Pr\"ufungszwecken vorgelegt und keine anderen als die
	angegebenen Hilfsmittel verwendet habe. S\"amtliche wissentlich verwendete
	Textausschnitte, Zitate oder Inhalte anderer Verfasser wurden ausdr\"ucklich als
	solche gekennzeichnet.
	
	Bremen, den \makeatletter\@date\makeatother
	
	\vspace*{1em}
	\rule{15em}{0.16667pt}\\
	\author{Naser Azizi, Sorin Arion}
	\makeatletter\@author\makeatother
	
	
	\normalem
	%Abstrakt
	%\addcontentsline{toc}{chapter}{Abstract}


	\chapter*{Motivation}
	
	Although industrial robots have been considered standard for many years, the widespread adoption of AI-based autonomous household robots is still far from being anticipated in every household.
	Obstacles such as unfamiliar surroundings and a limited knowledge base can significantly hinder the ability of an autonomous robot to effectively respond in a given situation, hence, the robot must precisely plan for every minor detail, even those that might appear trivial to us humans.
	To achieve the defined objectives, a proficient and comprehensive robotic system including Control, Perception, Navigation, and Knowledge is necessary.
	An illustrative goal could be cooking. To execute this task, the robot must possess various capabilities, including placing items, grasping objects, pouring liquids, cutting ingredients, and mixing components.	Those things seem trivial to us humans, but for the robots it requires a lot of Implementation and information.
	In this thesis, we aim to introduce an approach detailing how a robotic system can execute mixing actions effectively.	
	The challenges arise from the diverse methods of mixing, which further vary based on the specific ingredients being combined. Additionally, the consideration of compatible Mixing Tools and Mixing Containers is crucial, as not every tool can be utilized with every container.
	Through the incorporation of a knowledge graph containing rules related to to various actions, we aim to empower the robotic system to make informed decisions on the appropriate motions to employ. This decision-making process will take into account the specific task at hand and the involved ingredients.
	By incorporating this knowledge graph, we aim to advance towards autonomous robots capable of engaging in cooking activities.

	\chapter*{Introduction}
	First, we intend to analyze existing research, particularly examining studies that delve into diverse mixing motions and exploring related systems that contain a queryable knowledge graph.
	In the following section, we will introduce the reader with the frameworks employed to accomplish our objectives.	
	Subsequently, we will present our approach to data acquisition and data representation.
	In this chapter, we aim to explain the mixing tasks under consideration and articulate our methodology for representing this data to enable effective querying.
	In the subsequent section, we will illustrate the implementation of our rules, with the ultimate aim of determining the appropriate motion to be employed.
	To validate our concept, we have opted to simulate various scenarios and assess their outcomes, demonstrating the efficacy of our implemented system.
	Finally, we will delve into the discussion of our results and draw conclusions to wrap up our work.
	\chapter*{Related Work}
	\begin{itemize}
		\item https://fluidlab2023.github.io/
		\item 
	\end{itemize}

	\chapter*{How does the robot works / plans, parameters, used tools, etc}
	\begin{itemize}
		\item PyCram
		\item Pr2
		\item Ontology
		\item OWLReady
		\item SWRL
		\item 
	\end{itemize}
	\chapter*{Data acquisition}
	\begin{itemize}
		\item WikiHow Hyponyms
		\item Video analysis
		\item Structured the required parameters for the knowledge base
	\end{itemize}
	\chapter*{Data representation}
	\begin{itemize}
		\item Ingredients, Tools, etc.
		\item Different Tasks/Actions maps on different Motions
		\item Motion dependencies
		\item Fobidden Tools with Containers
		\item How is an action defined
		\item Examples of using the knowledge graph
	\end{itemize}
	\chapter*{Implementation}
	\begin{itemize}
		\item OWLReady
		\item SWRL
		\item Inferring parameters for the actual plan
	\end{itemize}
	\chapter*{Simulation}
	\begin{itemize}
		\item BulletWorld as Simulation platform -> how is the robot actually shown, which limitations exist
		\item Quering the knowledge base
		\item Inferring the parameters in the Simulation
		\item Task execution with multiple sequential Motions
	\end{itemize}
	\section*{Simulation to RealWorld gap}
	\begin{itemize}
		\item Big Problem: Uncertaininty
		\item Perception Solution as first approach: RoboKudo
		\item Data acquisition: Blenderproc
		\item Model Training: YoloV8
		\item Results showing the real world Perception.
	\end{itemize}
	\chapter*{Evaluation}
	\begin{itemize}
		\item Evaluation of multiple motions multiple times
		\item Results of that
	\end{itemize}
	\chapter*{Summary / Fazit}
avc
\end{document}

